% !TEX TS-program = pdflatex
% !TEX encoding = UTF-8 Unicode

\documentclass{beamer}

\mode<presentation>
{
  \usetheme{Copenhagen}
  \usecolortheme{beaver}
  % or ...

  \setbeamercovered{transparent}
  % or whatever (possibly just delete it)
}
\newcommand{\pGraphics}{../local/graphics}

\usepackage[english]{babel}
\usepackage[utf8]{inputenc}
\usepackage[]{hyperref}

\usepackage{times}
\usepackage[T1]{fontenc}
\usepackage{overpic}
% Or whatever. Note that the encoding and the font should match. If T1
% does not look nice, try deleting the line with the fontenc.
\begin{document}

\title[Confidence Assessment for IIR\#2] % (optional, use only with long paper titles)
{Confidence Assessment for IIR\#2}

\subtitle
{Presentation Subtitle} % (optional)

\author[Teague, Topa] % (optional, use only with lots of authors)
{F.~Author\inst{1} \and S.~Another\inst{2}}
% - Use the \inst{?} command only if the authors have different
%   affiliation.

\institute[Universities of Somewhere and Elsewhere] % (optional, but mostly needed)
{
  \inst{1}%
  Department of Computer Science\\
  University of Somewhere
  \and
  \inst{2}%
  Department of Theoretical Philosophy\\
  University of Elsewhere}
% - Use the \inst command only if there are several affiliations.
% - Keep it simple, no one is interested in your street address.

\date[Short Occasion] % (optional)
{Date / Occasion}

\subject{Talks}
%
\begin{frame}{Confidence Plots: Top Line}
    \centering
    \footnotesize
    
    % Set up a customized table
    % - First column (r): Right-aligns the numbering to line up the rightmost decimals
    % - @{\hspace{0.3cm}}: Forces a clean, consistent structural indent gap
    % - Second column (l): Left-aligns your system text entries
    
    \begin{tabular}{ l @{\hspace{0.3cm}} l c }
        \textbf{WBS} & \textbf{Product Element} & \textbf{Confidence Scores}  \\ \hline
       	\noalign{\vskip 2mm}
         \raisebox{0.8cm}{1.0}   
         & \raisebox{0.8cm}{Strategic Missile Weapon System} 
         & \includegraphics[ scale = 0.4 ]{\pGraphics/img-conf-overall-01.pdf} \\
    \end{tabular}
\end{frame}

\begin{frame}{\bf{Confidence Numbers: Top Line}}
	\center
	\begin{overpic}[ width = 0.55\textwidth ]{\pGraphics/ellipse-05.pdf}
	\end{overpic}
\end{frame}
%
\begin{frame}{\bf{Confidence Numbers: Top Line}}
\begin{table}[htbp]
\centering
\begin{tabular}{l rrr rrr r }
&  \multicolumn{3}{c}{Overall} &  \multicolumn{3}{c}{Elemental} \\
WBS 
& \multicolumn{1}{c}{$\mu$} 
&  \multicolumn{1}{c}{$\sigma$} 
& \multicolumn{1}{c}{$m$} 
&  \multicolumn{1}{c}{$\mu$}
&  \multicolumn{1}{c}{$\sigma$}
&  \multicolumn{1}{c}{$m$}
&  \multicolumn{1}{c}{$\mu_{O} - \mu_{E}$}  \\ \hline
       	\noalign{\vskip 1mm}
1.0         & \color{blue}{$0.23$}          & \color{purple}{$0.43$}
	          & $42$            & \color{blue}{$-0.15$}          & \color{purple}{$0.30$}          & $41$     & $0.38$        \\
\end{tabular}
%\caption{Sample data table matching your column types.}
\label{tab:test_data}
\end{table}
\end{frame}
%
\begin{frame}{\bf{Confidence Numbers: Top Line}}
	\begin{table}[htbp]
		\centering
		\begin{tabular}{l rr rr r }
			&  \multicolumn{2}{c}{Overall} &  \multicolumn{2}{c}{Elemental} \\
			WBS 
			& \multicolumn{1}{c}{$\mu \pm \sigma$} 
			& \multicolumn{1}{c}{$m$} 
			&  \multicolumn{1}{c}{$\mu \pm \sigma$}
			&  \multicolumn{1}{c}{$m$}
			&  \multicolumn{1}{c}{$\mu_{O} - \mu_{E}$}  \\ \hline
			\noalign{\vskip 1mm}
			1.0         & $\color{blue}{0.23}   \pm \color{purple}{0.43}$
			& $42$      & $\color{blue}{-0.15} \pm \color{purple}{0.30}$          & $41$     & $0.38$        \\
		\end{tabular}
		%\caption{Sample data table matching your column types.}
		\label{tab:test_data}
	\end{table}
\end{frame}
%
\begin{frame}{\bf{WBS Breakout: Top Level}}
    \centering
    \footnotesize
    
    % Set up a customized table
    % - First column (r): Right-aligns the numbering to line up the rightmost decimals
    % - @{\hspace{0.3cm}}: Forces a clean, consistent structural indent gap
    % - Second column (l): Left-aligns your system text entries
    
    \begin{tabular}{ l @{\hspace{0.3cm}} l }
        \textbf{WBS} & \textbf{Product Element} \\\hline
         1.0   & Strategic Missile Weapon System \\
        1.1 & Missile System \\
        1.2 & Launch System \\
        1.3 & Command \& Control System \\
        1.4 & Ground Systems \\
        1.5 & System Integration \\
        1.6 & System Test and Evaluation \\
        1.7 & Peculiar Support Equipment \\
        1.8 & Training
    \end{tabular}

\end{frame}
%
\begin{frame}{\bf{WBS 1.2 Breakout}}
    \centering
    \footnotesize
    
    \begin{tabular}{ l @{\hspace{0.3cm}} l }
        \textbf{WBS} & \textbf{Product Element} \\
         1.2   & Launch System \\
        1.2.1 & Launch Facility Equipment \\
        1.2.2 & Launch Support Equipment \\
        1.2.3 & Launch Control Interfaces \\
        1.2.4 & Environmental Control Equipment \\
    \end{tabular}
\end{frame}

\section*{Summary}

\begin{frame}{Summary}

  % Keep the summary *very short*.
  \begin{itemize}
  \item
    The \alert{first main message} of your talk in one or two lines.
  \item
    The \alert{second main message} of your talk in one or two lines.
  \item
    Perhaps a \alert{third message}, but not more than that.
  \end{itemize}
  
  % The following outlook is optional.
  \vskip0pt plus.5fill
  \begin{itemize}
  \item
    Outlook
    \begin{itemize}
    \item
      Something you haven't solved.
    \item
      Something else you haven't solved.
    \end{itemize}
  \end{itemize}
\end{frame}


\end{document}


